\source{Лекции, с. 55–60; Карпов, §4.1–4.2; детали лекционного доказательства Ловаса восстановлены}
\begin{definition}[Раскраска]
\(\chi(G)\) — минимальное число цветов правильной раскраски вершин;
\(\Delta=\Delta(G)\) — максимальная степень. В этом билете граф простой и связный.
Жадный алгоритм всегда даёт \(\chi(G)\le\Delta+1\): в любом порядке
каждая очередная вершина видит не более \(\Delta\) уже запрещённых цветов.
\end{definition}
Граф называется \textbf{регулярным}, если все его вершины имеют одну степень.
Если эта степень равна \(d\), граф называют \(d\)-регулярным.\par\medskip
\begin{lemma}[Вершина неполной степени]\label{lem:greedy}
Если \(G\) связен, \(\Delta(G)\le d\), \(d\ge1\), и есть вершина \(r\)
степени меньше \(d\), то \(\chi(G)\le d\).
\end{lemma}
\begin{proof}
Выберем остовное дерево с корнем \(r\) и красим вершины в порядке
убывания расстояния до корня (порядок на одном уровне произволен).
У каждой \(x\ne r\) родитель ещё не окрашен, поэтому окрашенных
соседей не больше \(d-1\). Из \(d\) цветов всегда есть свободный.
Корень красим последним; у него всего меньше \(d\) соседей,
так что свободный цвет тоже найдётся.
\end{proof}
\begin{lemma}[Тройка Ловаса]\label{lem:lovasz}
Если граф двусвязен, неполон и \(\delta(G)\ge3\), существуют вершины
\(a,b,c\), для которых \(ab,bc\in E\), \(ac\notin E\),
а \(G-\{a,c\}\) связен.
\end{lemma}
\begin{proof}
Если граф трёхсвязен (имеет не менее четырёх вершин и остаётся связным
после удаления любых одной или двух), возьмём кратчайший путь между
двумя несмежными вершинами. Его первые три вершины \(a,b,c\)
удовлетворяют \(ac\notin E\), иначе путь не кратчайший.
Связность после удаления \(a,c\) следует из трёхсвязности.

Если граф не трёхсвязен, существует разделяющая пара \(\{b,z\}\).
Граф \(G'=G-b\) связен по двусвязности \(G\), но имеет точку сочленения \(z\).
По дереву блоков (билет 5) у \(G'\) есть два разных крайних блока
\(B_a,B_c\). Обозначим единственные точки их присоединения к остальному
графу через \(s_a,s_c\) (они могут совпасть).
Вершина \(b\) смежна с некоторой внутренней вершиной каждого крайнего
блока. Иначе удаление \(s_a\), соответственно \(s_c\), отделяло бы
внутренние вершины блока и в исходном графе \(G\), вопреки двусвязности.
Выберем таких соседей \(a\in B_a\setminus\{s_a\}\),
\(c\in B_c\setminus\{s_c\}\).
Ребра \(ac\) нет, поскольку внутренние вершины разных крайних блоков
не соединены рёбрами.

Крайний блок не может состоять из одного ребра: его внутренняя вершина
имела бы степень 1 в \(G'\) и не больше 2 в \(G\), вопреки \(\delta(G)\ge3\).
Следовательно, оба выбранных блока невырожденно двусвязны.
Удаление \(a\) из \(B_a\) и \(c\) из \(B_c\) не нарушает их связности;
точки присоединения остаются, и потому \(G'-\{a,c\}\) связен.
У \(b\) не менее трёх соседей, так что после удаления двух вершин
хотя бы один сосед остаётся. Добавление \(b\) сохраняет связность.
\end{proof}
\begin{theorem}[Брукс]
Если связный простой граф не является ни полным графом,
ни циклом нечётной длины, то
\[
 \boxed{\chi(G)\le\Delta(G).}
\]
Исключения удовлетворяют \(\chi(K_n)=n=\Delta+1\),
\(\chi(C_{2r+1})=3=\Delta+1\).
\end{theorem}
\begin{proof}
При \(\Delta\le1\) связный граф — \(K_1\) или \(K_2\), оба исключены.
При \(\Delta=2\) связный граф — путь либо цикл:
если есть вершина степени 1, двигаясь из неё по единственному продолжению,
проходим весь граф, иначе возникла бы несвязность; если все степени 2,
тот же обход замыкается в цикл и исчерпывает граф.
Путь и чётный цикл красятся в два цвета поочерёдно;
нечётный цикл исключён.

Пусть \(d=\Delta\ge3\). Если хотя бы одна вершина имеет степень меньше
\(d\), применяем лемму \ref{lem:greedy}. Остаётся \(d\)-регулярный граф.
Если есть точка сочленения \(v\), обозначим компоненты \(G-v\) через
\(H_1,\ldots,H_t\), \(t\ge2\). Каждый граф \(G_i=G[V(H_i)\cup\{v\}]\)
связен, имеет максимальную степень не больше \(d\), а степень \(v\)
в нём меньше \(d\), поскольку хотя бы один сосед \(v\) находится
в другой компоненте. По первой лемме каждый \(G_i\) красится
в \(d\) цветов. Перестановкой названий цветов делаем цвет \(v\)
одинаковым во всех раскрасках и объединяем их. Между разными \(H_i\)
рёбер нет, поэтому раскраска правильна.

Наконец, пусть \(G\) двусвязен. По лемме \ref{lem:lovasz} есть
\(a,b,c\) с несмежными \(a,c\), общим соседом \(b\) и связным
\(G-\{a,c\}\). Красим \(a,c\) в один цвет. В оставшемся графе
берём остовное дерево с корнем \(b\) и красим снизу вверх.
У каждой вершины, кроме \(b\), остаётся неокрашенный родитель,
поэтому запрещено не более \(d-1\) цветов, и выбор возможен.
Корень \(b\) видит \(d\) соседей, но два из них, \(a,c\), имеют
одинаковый цвет. Следовательно, занято не более \(d-1\) цветов;
свободным цветом красим \(b\). Получена правильная \(d\)-раскраска.

В полном графе все вершины попарно смежны, поэтому нужны \(n\) цветов.
Нечётный цикл нельзя раскрасить в два цвета из-за чередования цветов
вдоль обхода; трёх цветов достаточно, если одну вершину покрасить
третьим цветом, а оставшийся путь — двумя. Исключения проверены.
\end{proof}
Для несвязного графа утверждение применяют к каждой компоненте отдельно;
хроматическое число всего графа — максимум чисел его компонент, поскольку
между ними нет ограничений на повторное использование цветов.
